
Douglas Richard Hofstadter (born February 15, 1945 -) is an American
scholar of cognitive science, physics, and comparative literature whose
research includes concepts such as the sense of self in relation to the
external world, consciousness, analogy-making, artistic creation, literary
translation, and discovery in mathematics and physics. His 1979 book
Gödel, Escher, Bach: An Eternal Golden Braid won both the Pulitzer Prize
for general nonfiction and a National Book Award (at that time called The
American Book Award) for Science. His 2007 book I Am a Strange Loop
won the Los Angeles Times Book Prize for Science and Technology.
Hofstadter was born in New York City to Jewish parents: Nobel Prize-
winning physicist Robert Hofstadter and Nancy Givan Hofstadter. He
grew up on the campus of Stanford University, where his father was a
professor, and attended the International School of Geneva in 1958–59. He
graduated with Distinction in mathematics from Stanford University in
1965, and received his Ph.D. in physics from the University of Oregon in
1975, where his study of the energy levels of Bloch electrons in a magnetic
field led to his discovery of the fractal known as Hofstadter's butterfly.




                                    Contents
List of Illustrations                                         xiv
Words of Thanks                                               xix

Part I: GEB
        Introduction: A Musico-Logical Offering               3
               Three-Part Invention                           29
        Chapter I: The MU-puzzle                              33
               Two-Part Invention                             43
        Chapter II: Meaning and Form in Mathematics           46
               Sonata for Unaccompanied Achilles              61
        Chapter III: Figure and Ground                        64
               Contracrostipunctus                            75
        Chapter IV: Consistency, Completeness, and Geometry   82
               Little Harmonic Labyrinth                      103
        Chapter V: Recursive Structures and Processes         127
               Canon by Intervallic Augmentation              153
        Chapter VI: The Location of Meaning                   158
               Chromatic Fantasy, And Feud                    177
        Chapter VII: The Propositional Calculus               181
               Crab Canon                                     199
        Chapter VIII: Typographical Number Theory             204
               A Mu Offering                                  231
        Chapter IX: Mumon and Gödel                           246





Part II EGB

      Prelude ...                                          275
Chapter X: Levels of Description, and Computer Systems     285
      Ant Fugue                                            311
Chapter XI: Brains and Thoughts                            337
      English French German Suit                           366
Chapter XII: Minds and Thoughts                            369
      Aria with Diverse Variations                         391
Chapter XIII: BlooP and FlooP and GlooP                    406
      Air on G's String                                    431
Chapter XIV: On Formally Undecidable Propositions of TNT
   and Related Systems                                     438
      Birthday Cantatatata ...                             461
Chapter XV: Jumping out of the System                      465
      Edifying Thoughts of a Tobacco Smoker                480
Chapter XVI: Self-Ref and Self-Rep                         495
      The Magn fierab, Indeed                              549
Chapter XVII: Church, Turing, Tarski, and Others           559
      SHRDLU, Toy of Man's Designing                       586
Chapter XVIII: Artificial Intelligence: Retrospects        594
      Contrafactus                                         633
Chapter XIX: Artificial Intelligence: Prospects            641
      Sloth Canon                                          681
Chapter XX: Strange Loops, Or Tangled Hierarchies          684
      Six-Part Ricercar                                    720


Notes                                                      743
Bibliography                                               746
Credits                                                    757
Index                                                      759





                                             Overview
                                           Part I: GEB
Introduction: A Musico-Logical Offering. The book opens with the story of Bach's Musical
   Offering. Bach made an impromptu visit to King Frederick the Great of Prussia, and was
   requested to improvise upon a theme presented by the King. His improvisations formed the basis
   of that great work. The Musical Offering and its story form a theme upon which I "improvise"
   throughout the book, thus making a sort of "Metamusical Offering". Self-reference and the
   interplay between different levels in Bach are discussed: this leads to a discussion of parallel
   ideas in Escher's drawings and then Gödel’s Theorem. A brief presentation of the history of logic
   and paradoxes is given as background for Gödel’s Theorem. This leads to mechanical reasoning
   and computers, and the debate about whether Artificial Intelligence is possible. I close with an
   explanation of the origins of the book-particularly the why and wherefore of the Dialogues.

Three-Part Invention. Bach wrote fifteen three-part inventions. In this three-part Dialogue, the
   Tortoise and Achilles-the main fictional protagonists in the Dialogues-are "invented" by Zeno (as
   in fact they were, to illustrate Zeno's paradoxes of motion). Very short, it simply gives the flavor
   of the Dialogues to come.

Chapter I: The MU-puzzle. A simple formal system (the MIL'-system) is presented, and the reader
  is urged to work out a puzzle to gain familiarity with formal systems in general. A number of
  fundamental notions are introduced: string, theorem, axiom, rule of inference, derivation, formal
  system, decision procedure, working inside/outside the system.

Two-Part Invention. Bach also wrote fifteen two-part inventions. This two-part Dialogue was written
  not by me, but by Lewis Carroll in 1895. Carroll borrowed Achilles and the Tortoise from Zeno,
  and I in turn borrowed them from Carroll. The topic is the relation between reasoning, reasoning
  about reasoning, reasoning about reasoning about reasoning, and so on. It parallels, in a way,
  Zeno's paradoxes about the impossibility of motion, seeming to show, by using infinite regress,
  that reasoning is impossible. It is a beautiful paradox, and is referred to several times later in the
  book.

Chapter II: Meaning and Form in Mathematics. A new formal system (the pq-system) is
  presented, even simpler than the MIU-system of Chapter I. Apparently meaningless at first, its
  symbols are suddenly revealed to possess meaning by virtue of the form of the theorems they
  appear in. This revelation is the first important insight into meaning: its deep connection to
  isomorphism. Various issues related to meaning are then discussed, such as truth, proof, symbol
  manipulation, and the elusive concept, "form".

Sonata for Unaccompanied Achilles. A Dialogue which imitates the Bach Sonatas for
   unaccompanied violin. In particular, Achilles is the only speaker, since it is a transcript of one
   end of a telephone call, at the far end of which is the Tortoise. Their conversation concerns the
   concepts of "figure" and "ground" in various





   contexts- e.g., Escher's art. The Dialogue itself forms an example of the distinction, since
   Achilles' lines form a "figure", and the Tortoise's lines-implicit in Achilles' lines-form a "ground".

Chapter III: Figure and Ground. The distinction between figure and ground in art is compared to
  the distinction between theorems and nontheorems in formal systems. The question "Does a
  figure necessarily contain the same information as its ground\%" leads to the distinction between
  recursively enumerable sets and recursive sets.

Contracrostipunctus. This Dialogue is central to the book, for it contains a set of paraphrases of
  Gödel’s self-referential construction and of his Incompleteness Theorem. One of the paraphrases
  of the Theorem says, "For each record player there is a record which it cannot play." The
  Dialogue's title is a cross between the word "acrostic" and the word "contrapunctus", a Latin word
  which Bach used to denote the many fugues and canons making up his Art of the Fugue. Some
  explicit references to the Art of the Fugue are made. The Dialogue itself conceals some acrostic
  tricks.

Chapter IV: Consistency, Completeness, and Geometry. The preceding Dialogue is explicated to
  the extent it is possible at this stage. This leads back to the question of how and when symbols in
  a formal system acquire meaning. The history of Euclidean and non-Euclidean geometry is given,
  as an illustration of the elusive notion of "undefined terms". This leads to ideas about the
  consistency of different and possibly "rival" geometries. Through this discussion the notion of
  undefined terms is clarified, and the relation of undefined terms to perception and thought
  processes is considered.

Little Harmonic Labyrinth. This is based on the Bach organ piece by the same name. It is a playful
    introduction to the notion of recursive-i.e., nested structures. It contains stories within stories. The
    frame story, instead of finishing as expected, is left open, so the reader is left dangling without
    resolution. One nested story concerns modulation in music-particularly an organ piece which
    ends in the wrong key, leaving the listener dangling without resolution.

Chapter V: Recursive Structures and Processes. The idea of recursion is presented in many
  different contexts: musical patterns, linguistic patterns, geometric structures, mathematical
  functions, physical theories, computer programs, and others.

Canon by Intervallic Augmentation. Achilles and the Tortoise try to resolve the question, "Which
  contains more information-a record, or the phonograph which plays it This odd question arises
  when the Tortoise describes a single record which, when played on a set of different
  phonographs, produces two quite different melodies: B-A-C-H and C-A-G-E. It turns out,
  however, that these melodies are "the same", in a peculiar sense.

Chapter VI: The Location of Meaning. A broad discussion of how meaning is split among coded
  message, decoder, and receiver. Examples presented include strands of DNA, undeciphered
  inscriptions on ancient tablets, and phonograph records sailing out in space. The relationship of
  intelligence to "absolute" meaning is postulated.

Chromatic Fantasy, And Feud. A short Dialogue bearing hardly any resemblance, except in title, to
  Bach's Chromatic Fantasy and Fugue. It concerns the proper way to manipulate sentences so as
  to preserve truth-and in particular the question





   of whether there exist rules for the usage of the word "arid". This Dialogue has much in common
   with the Dialogue by Lewis Carroll.

Chapter VII: The Propositional Calculus. It is suggested how words such as .,and" can be
  governed by formal rules. Once again, the ideas of isomorphism and automatic acquisition of
  meaning by symbols in such a system are brought up. All the examples in this Chapter,
  incidentally, are "Zentences"-sentences taken from Zen koans. This is purposefully done,
  somewhat tongue-in-cheek, since Zen koans are deliberately illogical stories.

Crab Canon. A Dialogue based on a piece by the same name from the Musical Offering. Both are so
  named because crabs (supposedly) walk backwards. The Crab makes his first appearance in this
  Dialogue. It is perhaps the densest Dialogue in the book in terms of formal trickery and level-
  play. Gödel, Escher, and Bach are deeply intertwined in this very short Dialogue.

Chapter VIII: Typographical Number Theory. An extension of the Propositional Calculus called
  "TNT" is presented. In TNT, number-theoretical reasoning can be done by rigid symbol
  manipulation. Differences between formal reasoning and human thought are considered.

A Mu Offering. This Dialogue foreshadows several new topics in the book. Ostensibly concerned
  with Zen Buddhism and koans, it is actually a thinly veiled discussion of theoremhood and
  nontheoremhood, truth and falsity, of strings in number theory. There are fleeting references to
  molecular biology-particular) the Genetic Code. There is no close affinity to the Musical
  Offering, other than in the title and the playing of self-referential games.

Chapter IX: Mumon and Gödel. An attempt is made to talk about the strange ideas of Zen
  Buddhism. The Zen monk Mumon, who gave well known commentaries on many koans, is a
  central figure. In a way, Zen ideas bear a metaphorical resemblance to some contemporary ideas
  in the philosophy of mathematics. After this "Zennery", Gödel’s fundamental idea of Gödel-
  numbering is introduced, and a first pass through Gödel’s Theorem is made.

                                         Part II: EGB
Prelude ... This Dialogue attaches to the next one. They are based on preludes and fugues from
   Bach's Well-Tempered Clavier. Achilles and the Tortoise bring a present to the Crab, who has a
   guest: the Anteater. The present turns out to be a recording of the W.T.C.; it is immediately put
   on. As they listen to a prelude, they discuss the structure of preludes and fugues, which leads
   Achilles to ask how to hear a fugue: as a whole, or as a sum of parts? This is the debate between
   holism and reductionism, which is soon taken up in the Ant Fugue.

Chapter X: Levels of Description, and Computer Systems. Various levels of seeing pictures,
  chessboards, and computer systems are discussed. The last of these is then examined in detail.
  This involves describing machine languages, assembly languages, compiler languages, operating
  systems, and so forth. Then the discussion turns to composite systems of other types, such as
  sports teams, nuclei, atoms, the weather, and so forth. The question arises as to how man
  intermediate levels exist-or indeed whether any exist.





   …Ant Fugue. An imitation of a musical fugue: each voice enters with the same statement. The
   theme-holism versus reductionism-is introduced in a recursive picture composed of words
   composed of smaller words. etc. The words which appear on the four levels of this strange picture
   are "HOLISM", "REDLCTIONIsM", and "ML". The discussion veers off to a friend of the
   Anteater's Aunt Hillary, a conscious ant colony. The various levels of her thought processes are
   the topic of discussion. Many fugal tricks are ensconced in the Dialogue. As a hint to the reader,
   references are made to parallel tricks occurring in the fugue on the record to which the foursome
   is listening. At the end of the Ant Fugue, themes from the Prelude return. transformed
   considerably.

Chapter XI: Brains and Thoughts. "How can thoughts he supported by the hardware of the brain is
  the topic of the Chapter. An overview of the large scale and small-scale structure of the brain is
  first given. Then the relation between concepts and neural activity is speculatively discussed in
  some detail.

English French German Suite. An interlude consisting of Lewis Carroll's nonsense poem
  "Jabberwocky`' together with two translations: one into French and one into German, both done
  last century.

Chapter XII: Minds and Thoughts. The preceding poems bring up in a forceful way the question
  of whether languages, or indeed minds, can be "mapped" onto each other. How is communication
  possible between two separate physical brains: What do all human brains have in common? A
  geographical analogy is used to suggest an answer. The question arises, "Can a brain be
  understood, in some objective sense, by an outsider?"

Aria with Diverse Variations. A Dialogue whose form is based on Bach's Goldberg Variations, and
   whose content is related to number-theoretical problems such as the Goldbach conjecture. This
   hybrid has as its main purpose to show how number theory's subtlety stems from the fact that
   there are many diverse variations on the theme of searching through an infinite space. Some of
   them lead to infinite searches, some of them lead to finite searches, while some others hover in
   between.

Chapter XIII: BlooP and FlooP and GlooP. These are the names of three computer languages.
  BlooP programs can carry out only predictably finite searches, while FlooP programs can carry
  out unpredictable or even infinite searches. The purpose of this Chapter is to give an intuition for
  the notions of primitive recursive and general recursive functions in number theory, for they are
  essential in Gödel’s proof.

Air on G's String. A Dialogue in which Gödel’s self-referential construction is mirrored in words.
   The idea is due to W. V. O. Quine. This Dialogue serves as a prototype for the next Chapter.

Chapter XIV: On Formally Undecidable Propositions of TNT and Related Systems. This
  Chapter's title is an adaptation of the title of Gödel’s 1931 article, in which his Incompleteness
  Theorem was first published. The two major parts of Gödel’s proof are gone through carefully. It
  is shown how the assumption of consistency of TNT forces one to conclude that TNT (or any
  similar system) is incomplete. Relations to Euclidean and non-Euclidean geometry are discussed.
  Implications for the philosophy of mathematics are gone into with some care.





Birthday Cantatatata ... In which Achilles cannot convince the wily and skeptical Tortoise that today
   is his (Achilles') birthday. His repeated but unsuccessful tries to do so foreshadow the
   repeatability of the Gödel argument.

Chapter XV: Jumping out of the System. The repeatability of Gödel’s argument is shown, with
  the implication that TNT is not only incomplete, but "essentially incomplete The fairly notorious
  argument by J. R. Lucas, to the effect that Gödel’s Theorem demonstrates that human thought
  cannot in any sense be "mechanical", is analyzed and found to be wanting.

Edifying Thoughts of a Tobacco Smoker. A Dialogue treating of many topics, with the thrust being
   problems connected with self-replication and self-reference. Television cameras filming
   television screens, and viruses and other subcellular entities which assemble themselves, are
   among the examples used. The title comes from a poem by J. S. Bach himself, which enters in a
   peculiar way.

Chapter XVI: Self-Ref and Self-Rep. This Chapter is about the connection between self-reference
  in its various guises, and self-reproducing entities e.g., computer programs or DNA molecules).
  The relations between a self-reproducing entity and the mechanisms external to it which aid it in
  reproducing itself (e.g., a computer or proteins) are discussed-particularly the fuzziness of the
  distinction. How information travels between various levels of such systems is the central topic of
  this Chapter.

The Magnificrab, Indeed. The title is a pun on Bach's Magnifacat in D. The tale is about the Crab,
   who gives the appearance of having a magical power of distinguishing between true and false
   statements of number theory by reading them as musical pieces, playing them on his flute, and
   determining whether they are "beautiful" or not.

Chapter XVII: Church, Turing, Tarski, and Others. The fictional Crab of the preceding Dialogue
  is replaced by various real people with amazing mathematical abilities. The Church-Turing
  Thesis, which relates mental activity to computation, is presented in several versions of differing
  strengths. All are analyzed, particularly in terms of their implications for simulating human
  thought mechanically, or programming into a machine an ability to sense or create beauty. The
  connection between brain activity and computation brings up some other topics: the halting
  problem of Turing, and Tarski's Truth Theorem.

SHRDLU, Toy of Man's Designing. This Dialogue is lifted out of an article by Terry Winograd on
  his program SHRDLU: only a few names have been changed. In it. a program communicates
  with a person about the so-called "blocks world" in rather impressive English. The computer
  program appears to exhibit some real understanding-in its limited world. The Dialogue's title is
  based on Jesu, joy of Mans Desiring, one movement of Bach's Cantata 147.

Chapter XVIII: Artificial Intelligence: Retrospects, This Chapter opens with a discussion of the
  famous "Turing test"-a proposal by the computer pioneer Alan Turing for a way to detect the
  presence or absence of "thought" in a machine. From there, we go on to an abridged history of
  Artificial Intelligence. This covers programs that can-to some degree-play games, prove
  theorems, solve problems, compose music, do mathematics, and use "natural language" (e.g.,
  English).





Contrafactus. About how we unconsciously organize our thoughts so that we can imagine
  hypothetical variants on the real world all the time. Also about aberrant variants of this ability-
  such as possessed by the new character, the Sloth, an avid lover of French fries, and rabid hater of
  counterfactuals.

Chapter XIX: Artificial Intelligence: Prospects. The preceding Dialogue triggers a discussion of
  how knowledge is represented in layers of contexts. This leads to the modern Al idea of "frames".
  A frame-like way of handling a set of visual pattern puzzles is presented, for the purpose of
  concreteness. Then the deep issue of the interaction of concepts in general is discussed, which
  leads into some speculations on creativity. The Chapter concludes with a set of personal
  "Questions and Speculations" on Al and minds in general.

Sloth Canon. A canon which imitates a Bach canon in which one voice plays the same melody as
   another, only upside down and twice as slowly, while a third voice is free. Here, the Sloth utters
   the same lines as the Tortoise does, only negated (in a liberal sense of the term) and twice as
   slowly, while Achilles is free.

Chapter XX: Strange Loops, Or Tangled Hierarchies. A grand windup of many of the ideas
  about hierarchical systems and self-reference. It is concerned with the snarls which arise when
  systems turn back on themselves-for example, science probing science, government investigating
  governmental wrongdoing, art violating the rules of art, and finally, humans thinking about their
  own brains and minds. Does Gödel’s Theorem have anything to say about this last "snarl"? Are
  free will and the sensation of consciousness connected to Gödel’s Theorem? The Chapter ends by
  tying Gödel, Escher, and Bach together once again.

Six-Part Ricercar. This Dialogue is an exuberant game played with many of the ideas which have
   permeated the book. It is a reenactment of the story of the Musical Offering, which began the
   book; it is simultaneously a "translation" into words of the most complex piece in the Musical
   Offering: the Six-Part Ricercar. This duality imbues the Dialogue with more levels of meaning
   than any other in the book. Frederick the Great is replaced by the Crab, pianos by computers, and
   so on. Many surprises arise. The Dialogue's content concerns problems of mind, consciousness,
   free will, Artificial Intelligence, the Turing test, and so forth, which have been introduced earlier.
   It concludes with an implicit reference to the beginning of the book, thus making the book into
   one big self-referential loop, symbolizing at once Bach's music, Escher's drawings, and Gödel’s
   Theorem.





FIGURE 1. Johann Sebastian Bach, in 1748. From a painting by Elias Gottlieb
Hanssmann.




Introduction: A Musico-Logical Offering                                 10

                                    Introduction:

                            A Musico-Logical Offering
Author:
   FREDERICK THE GREAT, King of Prussia, came to power in 1740. Although he is
remembered in history books mostly for his military astuteness, he was also devoted to
the life of the mind and the spirit. His court in Potsdam was one of the great centers of
intellectual activity in Europe in the eighteenth century. The celebrated mathematician
Leonhard Euler spent twenty-five years there. Many other mathematicians and scientists
came, as well as philosophers-including Voltaire and La Mettrie, who wrote some of their
most influential works while there.
   But music was Frederick's real love. He was an avid flutist and composer. Some of his
compositions are occasionally performed even to this day. Frederick was one of the first
patrons of the arts to recognize the virtues of the newly developed "piano-forte" ("soft-
loud"). The piano had been developed in the first half of the eighteenth century as a
modification of the harpsichord. The problem with the harpsichord was that pieces could
only be played at a rather uniform loudness-there was no way to strike one note more
loudly than its neighbors. The "soft-loud", as its name implies, provided a remedy to this
problem. From Italy, where Bartolommeo Cristofori had made the first one, the soft-loud
idea had spread widely. Gottfried Silbermann, the foremost German organ builder of the
day, was endeavoring to make a "perfect" piano-forte. Undoubtedly King Frederick was
the greatest supporter of his efforts-it is said that the King owned as many as fifteen
Silbermann pianos!

                                          Bach

Frederick was an admirer not only of pianos, but also of an organist and composer by the
name of J. S. Bach. This Bach's compositions were somewhat notorious. Some called
them "turgid and confused", while others claimed they were incomparable masterpieces.
But no one disputed Bach's ability to improvise on the organ. In those days, being an
organist not only meant being able to play, but also to extemporize, and Bach was known
far and wide for his remarkable extemporizations. (For some delightful anecdotes about
Bach's extemporization, see The Bach Reader, by H. T. David and A. Mendel.)
   In 1747, Bach was sixty-two, and his fame, as well as one of his sons, had reached
Potsdam: in fact, Carl Philipp Emanuel Bach was the Capellmeister (choirmaster) at the
court of King Frederick. For years the King had let it be known, through gentle hints to
Philipp Emanuel, how




Introduction: A Musico-Logical Offering                                                11

pleased he would be to have the elder Bach come and pay him a visit; but this wish had
never been realized. Frederick was particularly eager for Bach to try out his new
Silbermann pianos, which lie (Frederick) correctly foresaw as the great new wave in
music.
   It was Frederick's custom to have evening concerts of chamber music in his court.
Often he himself would be the soloist in a concerto for flute Here we have reproduced a
painting of such an evening by the German painter Adolph von Menzel, who, in the
1800's, made a series of paintings illustrating the life of Frederick the Great. At the
cembalo is C. P. E. Bach, and the figure furthest to the right is Joachim Quantz, the
King's flute master-and the only person allowed to find fault with the King's flute
playing. One May evening in 1747, an unexpected guest showed up. Johann Nikolaus
Forkel, one of Bach's earliest biographers, tells the story
as follows:
 One evening, just as lie was getting his flute ready, and his musicians were ssembled,
 an officer brought him a list of the strangers who had arrived. With his flute in his hand
 he ran ever the list, but immediately turned to the assembled musicians, and said, with a
 kind of agitation, "Gentlemen, old Bach is come." The Hute was now laid aside, and old
 Bach, who had alighted at his son's lodgings, was immediately summoned to the Palace.
 Wilhelm Friedemann, who accompanied his father, told me this story, and I must say
 that 1 still think with pleasure on the manner in which lie related it. At that time it was
 the fashion to make rather prolix compliments. The first appearance of J. S. Bach before
 se great a King, who did not even give him time to change his traveling dress for a
 black chanter's gown, must necessarily be attended with many apologies. I will net here
 dwell en these apologies, but merely observe, that in Wilhelm Friedemann's mouth they
 made a formal Dialogue between the King and the Apologist.
 But what is mere important than this is that the King gave up his Concert for this
 evening, and invited Bach, then already called the Old Bach, to try his fortepianos,
 made by Silbermann, which steed in several rooms of the palace. [Forkel here inserts
 this footnote: "The pianofortes manufactured by Silbermann, of Frevberg, pleased the
 King se much, that he resolved to buy them all up. He collected fifteen. I hear that they
 all now stand unfit for use in various corners of the Royal Palace."] The musicians went
 with him from room to room, and Bach was invited everywhere to try them and to play
 unpremeditated compositions. After he had gene en for some time, he asked the King to
 give him a subject for a Fugue, in order to execute it immediately without any
 preparation. The King admired the learned manner in which his subject was thus
 executed extempore: and, probably to see hew far such art t could be carried, expressed
 a wish to hear a Fugue with six Obligato parts. But as it is not every subject that is fit
 for such full harmony, Bach chose one himself, and immediately executed it to the
 astonishment of all present in the same magnificent and learned manner as he had done
 that of the King. His Majesty desired also to hear his performance en the organ. The
 next day therefore Bach was taken to all the organs in Potsdam, as lie had before been
 to Silbermann's fortepianos. After his return to Leipzig, he composed the subject, which
 he had received from the King, in three and six parts. added several artificial passages
 in strict canon to it, and had it engraved, under the title of "Musikalisches Opfer"
 [Musical Offering], and dedicated it to the Inventor.'




Introduction: A Musico-Logical Offering                                                  12

Introduction: A Musico-Logical Offering   13

FIGURE 3. The Royal Theme.

   When Bach sent a copy of his Musical Offering to the King, he included a dedicatory
letter, which is of interest for its prose style if nothing else rather submissive and
flattersome. From a modern perspective it seems comical. Also, it probably gives
something of the flavor of Bach's apology for his appearance.2

  MOST GRACIOUS KING!

  In deepest humility I dedicate herewith to Your Majesty a musical offering, the
  noblest part of which derives from Your Majesty's own august hand. With awesome
  pleasure I still remember the very special Royal grace when, some time ago, during
  my visit in Potsdam, Your Majesty's Self deigned to play to me a theme for a fugue
  upon the clavier, and at the same time charged me most graciously to carry it out in
  Your Majesty's most august presence. To obey Your Majesty's command was my most
  humble dim. I noticed very soon, however, that, for lack of necessary preparation, the
  execution of the task did not fare as well as such an excellent theme demanded. I
  resoled therefore and promptly pledged myself to work out this right Royal theme
  more fully, and then make it known to the world. This resolve has now been carried
  out as well as possible, and it has none other than this irreproachable intent, to glorify,
  if only in a small point, the fame of a monarch whose greatness and power, as in all
  the sciences of war and peace, so especially in music, everyone must admire and
  revere. I make bold to add this most humble request: may Your Majesty deign to
  dignify the present modest labor with a gracious acceptance, and continue to grant
  Your Majesty's most august Royal grace to

                                             Your Majesty's
                                                   most humble and obedient servant,
                                                           THE AUTHOR
Leipzig, July 7 1747

  Some twenty-seven years later, when Bach had been dead for twentyfour years, a Baron
named Gottfried van Swieten-to whom, incidentally, Forkel dedicated his biography of
Bach, and Beethoven dedicated his First Symphony-had a conversation with King
Frederick, which he reported as follows:

  He [Frederick] spoke to me, among other things, of music, and of a great organist
  named Bach, who has been for a while in Berlin. This artist [Wilhelm Friedemann
  Bach] is endowed with a talent superior, in depth of harmonic knowledge and power
  of execution, to any 1 have heard or can imagine, while those who knew his father
  claim that he, in turn, was even greater. The King



Introduction: A Musico-Logical Offering                                                   14

   is of this opinion, and to prove it to me he sang aloud a chromatic fugue subject which
   he had given this old Bach, who on the spot had made of it a fugue in four parts, then
   in five parts, and finally in eight parts.'

   Of course there is no way of knowing whether it was King Frederick or Baron van
Swieten who magnified the story into larger-than-life proportions. But it shows how
powerful Bach's legend had become by that time. To give an idea of how extraordinary a
six-part fugue is, in the entire Well-Tempered Clavier by Bach, containing forty-eight
Preludes and Fugues, only two have as many as five parts, and nowhere is there a six-part
fugue! One could probably liken the task of improvising a six-part fugue to the playing of
sixty simultaneous blindfold games of chess, and winning them all. To improvise an
eight-part fugue is really beyond human capability.
   In the copy which Bach sent to King Frederick, on the page preceding the first sheet of
music, was the following inscription:




FIG URE 4.

("At the King's Command, the Song and the Remainder Resolved with Canonic Art.")
Here Bach is punning on the word "canonic", since it means not only "with canons" but
also "in the best possible way". The initials of this inscription are

                                    RICERCAR

-an Italian word, meaning "to seek". And certainly there is a great deal to seek in the
Musical Offering. It consists of one three-part fugue, one six-part fugue, ten canons, and a
trio sonata. Musical scholars have concluded that the three-part fugue must be, in
essence, identical with the one which Bach improvised for King Frederick. The six-part
fugue is one of Bach's most complex creations, and its theme is, of course, the Royal
Theme. That theme, shown in Figure 3, is a very complex one, rhythmically irregular and
highly chromatic (that is, filled with tones which do not belong to the key it is in). To
write a decent fugue of even two voices based on it would not be easy for the average
musician!
   Both of the fugues are inscribed "Ricercar", rather than "Fuga". This is another
meaning of the word; "ricercar" was, in fact, the original name for the musical form now
known as "fugue". By Bach's time, the word "fugue" (or fuga, in Latin and Italian) had
become standard, but the term "ricercar" had survived, and now designated an erudite
kind of fugue, perhaps too austerely intellectual for the common ear. A similar usage
survives in English today: the word "recherche" means, literally, "sought out", but carries
the same kind of implication, namely of esoteric or highbrow cleverness.
   The trio sonata forms a delightful relief from the austerity of the fugues and canons,
because it is very melodious and sweet, almost dance-



Introduction: A Musico-Logical Offering                                                  15

able. Nevertheless, it too is based largely on the King's theme, chromatic and austere as it
is. It is rather miraculous that Bach could use such a theme to make so pleasing an
interlude.
  The ten canons in the Musical Offering are among the most sophisticated canons Bach
ever wrote. However, curiously enough, Bach himself never wrote them out in full. This
was deliberate. They were posed as puzzles to King Frederick. It was a familiar musical
game of the day to give a single theme, together with some more or less tricky hints, and
to let the canon based on that theme be "discovered" by someone else. In order to know
how this is possible, you must understand a few facts about canons.

                                    Canons and Fugues

The idea of a canon is that one single theme is played against itself. This is done by
having "copies" of the theme played by the various participating voices. But there are
means' ways to do this. The most straightforward of all canons is the round, such as
"Three Blind Mice", "Row, Row, Row Your Boat", or " Frere Jacques". Here, the theme
enters in the first voice and, after a fixed time-delay, a "copy" of it enters, in precisely the
same key. After the same fixed time-delay in the second voice, the third voice enters
carrying the theme, and so on. Most themes will not harmonize with themselves in this
way. In order for a theme to work as a canon theme, each of its notes must be able to
serve in a dual (or triple, or quadruple) role: it must firstly be part of a melody, and
secondly it must be part of a harmonization of the same melody. When there are three
canonical voices, for instance, each note of the theme must act in two distinct harmonic
ways, as well as melodically. Thus, each note in a canon has more than one musical
meaning; the listener's ear and brain automatically figure out the appropriate meaning, by
referring to context.
  There are more complicated sorts of canons, of course. The first escalation in
complexity comes when the "copies" of the theme are staggered not only in time, but also
in pitch; thus, the first voice might sing the theme starting on C, and the second voice,
overlapping with the first voice, might sing the identical theme starting five notes higher,
on G. A third voice, starting on the D yet five notes higher, might overlap with the first
two, and so on. The next escalation in complexity comes when the speeds of' the different
voices are not equal; thus, the second voice might sing twice as quickly, or twice as
slowly, as the first voice. The former is called diminution, the latter augmentation (since
the theme seems to shrink or to expand).
  We are not yet done! The next stage of complexity in canon construction is to invert the
theme, which means to make a melody which jumps down wherever the original theme
jumps up, and by exactly the same number of semitones. This is a rather weird melodic
transformation, but when one has heard many themes inverted, it begins to seem quite
natural. Bach was especially fond of inversions, and they show up often in his work-and
the Musical Offering is no exception. (For a simple example of




Introduction: A Musico-Logical Offering                                                      16

inversion, try the tune "Good King Wenceslas". When the original and its inversion are
sung together, starting an octave apart and staggered with a time-delay of two beats, a
pleasing canon results.) Finally, the most esoteric of "copies" is the retrograde copy-
where the theme is played backwards in time. A canon which uses this trick is
affectionately known as a crab canon, because of the peculiarities of crab locomotion.
Bach included a crab canon in the Musical Offering, needless to say. Notice that every
type of "copy" preserves all the information in the original theme, in the sense that the
theme is fully recoverable from any of the copies. Such an information preserving
transformation is often called an isomorphism, and we will have much traffic with
isomorphisms in this book.
   Sometimes it is desirable to relax the tightness of the canon form. One way is to allow
slight departures from perfect copying, for the sake of more fluid harmony. Also, some
canons have "free" voices-voices which do not employ the canon's theme, but which
simply harmonize agreeably with the voices that are in canon with each other.
   Each of the canons in the Musical Offering has for its theme a different variant of the
King's Theme, and all the devices described above for making canons intricate are
exploited to the hilt; in fact, they are occasionally combined. Thus, one three-voice canon
is labeled "Canon per Augmentationem, contrario Motu"; its middle voice is free (in fact,
it sings the Royal Theme), while the other two dance canonically above and below it,
using the devices of augmentation and inversion. Another bears simply the cryptic label
"Quaerendo invenietis" ("By seeking, you will discover"). All of the canon puzzles have
been solved. The canonical solutions were given by one of Bach's pupils, Johann Philipp
Kirnberger. But one might still wonder whether there are more solutions to seek!
   I should also explain briefly what a fugue is. A fugue is like a canon, in that it is
usually based on one theme which gets played in different voices and different keys, and
occasionally at different speeds or upside down or backwards. However, the notion of
fugue is much less rigid than that of canon, and consequently it allows for more
emotional and artistic expression. The telltale sign of a fugue is the way it begins: with a
single voice singing its theme. When it is done, then a second voice enters, either five
scale-notes up, or four down. Meanwhile the first voice goes on, singing the
"countersubject": a secondary theme, chosen to provide rhythmic, harmonic, and melodic
contrasts to the subject. Each of the voices enters in turn, singing the theme, often to the
accompaniment of the countersubject in some other voice, with the remaining voices
doing whatever fanciful things entered the composer's mind. When all the voices have
"arrived", then there are no rules. There are, to be sure, standard kinds of things to do-but
not so standard that one can merely compose a fugue by formula. The two fugues in the
Musical Offering are outstanding examples of fugues that could never have been
"composed by formula". There is something much deeper in them than mere fugality.
   All in all, the Musical Offering represents one of Bach's supreme accomplishments in
counterpoint. It is itself one large intellectual fugue, in




Introduction: A Musico-Logical Offering                                                   17

which many ideas and forms have been woven together, and in which playful double
meanings and subtle allusions are commonplace. And it is a very beautiful creation of the
human intellect which we can appreciate forever. (The entire work is wonderfully
described in the book f. S. Bach's Musical Offering, by H. T. David.)

                                An Endlessly Rising Canon

There is one canon in the Musical Offering which is particularly unusual. Labeled simply
"Canon per Tonos", it has three voices. The uppermost voice sings a variant of the Royal
Theme, while underneath it, two voices provide a canonic harmonization based on a
second theme. The lower of this pair sings its theme in C minor (which is the key of the
canon as a whole), and the upper of the pair sings the same theme displaced upwards in
pitch by an interval of a fifth. What makes this canon different from any other, however,
is that when it concludes-or, rather, seems to conclude-it is no longer in the key of C
minor, but now is in D minor. Somehow Bach has contrived to modulate (change keys)
right under the listener's nose. And it is so constructed that this "ending" ties smoothly
onto the beginning again; thus one can repeat the process and return in the key of E, only
to join again to the beginning. These successive modulations lead the ear to increasingly
remote provinces of tonality, so that after several of them, one would expect to be
hopelessly far away from the starting key. And yet magically, after exactly six such
modulations, the original key of C minor has been restored! All the voices are exactly one
octave higher than they were at the beginning, and here the piece may be broken off in a
musically agreeable way. Such, one imagines, was Bach's intention; but Bach indubitably
also relished the implication that this process could go on ad infinitum, which is perhaps
why he wrote in the margin "As the modulation rises, so may the King's Glory." To
emphasize its potentially infinite aspect, I like to call this the "Endlessly Rising Canon".
   In this canon, Bach has given us our first example of the notion of Strange Loops. The
"Strange Loop" phenomenon occurs whenever, by moving upwards (or downwards)
through the levels of some hierarchical system, we unexpectedly find ourselves right
back where we started. (Here, the system is that of musical keys.) Sometimes I use the
term Tangled Hierarchy to describe a system in which a Strange Loop occurs. As we go
on, the theme of Strange Loops will recur again and again. Sometimes it will be hidden,
other times it will be out in the open; sometimes it will be right side up, other times it will
be upside down, or backwards. "Quaerendo invenietis" is my advice to the reader.


                                           Escher

To my mind, the most beautiful and powerful visual realizations of this notion of Strange
Loops exist in the work of the Dutch graphic artist M. C. Escher, who lived from 1902 to
1972. Escher was the creator of some of the




Introduction: A Musico-Logical Offering                                                     18

                FIGURE 5. Waterfall, by M. C. Escher (lithograph, 1961).

most intellectually stimulating drawings of all time. Many of them have their origin in
paradox, illusion, or double-meaning. Mathematicians were among the first admirers of
Escher's drawings, and this is understandable because they often are based on
mathematical principles of symmetry or pattern ... But there is much more to a typical
Escher drawing than just symmetry or pattern; there is often an underlying idea, realized
in artistic form. And in particular, the Strange Loop is one of the most recurrent themes in
Escher's work. Look, for example, at the lithograph Waterfall (Fig. 5), and compare its
six-step endlessly falling loop with the six-step endlessly rising loop of the "Canon per
Tonos". The similarity of vision is




Introduction: A Musico-Logical Offering                                                  19

      FIGURE 6. Ascending and Descending, by M. C. Escher (lithograph, 1960).




Introduction: A Musico-Logical Offering                                         20

remarkable. Bach and Escher are playing one single theme in two different "keys": music
and art.
   Escher realized Strange Loops in several different ways, and they can be arranged
according to the tightness of the loop. The lithograph Ascending and Descending (Fig. 6),
in which monks trudge forever in loops, is the loosest version, since it involves so many
steps before the starting point is regained. A tighter loop is contained in Waterfall, which,
as we already observed, involves only six discrete steps. You may be thinking that there
is some ambiguity in the notion of a single "step"-for instance, couldn't Ascending and
Descending be seen just as easily as having four levels (staircases) as forty-five levels
(stairs)\% It is indeed true that there is an inherent




   FIGURE 7. Hand with Reflecting Globe. Self-portrait In, M. C. Escher (lithograph,
                                     1935).




Introduction: A Musico-Logical Offering                                                   21

Introduction: A Musico-Logical Offering   22

haziness in level-counting, not only in Escher pictures, but in hierarchical, many-level
systems. We will sharpen our understanding of this haziness later on. But let us not get
too distracted now' As we tighten our loop, we come to the remarkable Drawing Hands
(Fig. 135), in which each of two hands draws the other: a two-step Strange Loop. And
finally, the tightest of all Strange Loops is realized in Print Gallery (Fig. 142): a picture
of a picture which contains itself. Or is it a picture of a gallery which contains itself? Or
of a town which contains itself? Or a young man who contains himself'? (Incidentally, the
illusion underlying Ascending and Descending and Waterfall was not invented by Escher,
but by Roger Penrose, a British mathematician, in 1958. However, the theme of the
Strange Loop was already present in Escher's work in 1948, the year he drew Drawing
Hands. Print Gallery dates from 1956.)
   Implicit in the concept of Strange Loops is the concept of infinity, since what else is a
loop but a way of representing an endless process in a finite way? And infinity plays a
large role n many of Escher's drawings. Copies of one single theme often fit into each'
other, forming visual analogues to the canons of Bach. Several such patterns can be seen
in Escher's famous print Metamorphosis (Fig. 8). It is a little like the "Endlessly Rising
Canon": wandering further and further from its starting point, it suddenly is back. In the
tiled planes of Metamorphosis and other pictures, there are already suggestions of
infinity. But wilder visions of infinity appear in other drawings by Escher. In some of his
drawings, one single theme can appear on different levels of reality. For instance, one
level in a drawing might clearly be recognizable as representing fantasy or imagination;
another level would be recognizable as reality. These two levels might be the only
explicitly portrayed levels. But the mere presence of these two levels invites the viewer to
look upon himself as part of yet another level; and by taking that step, the viewer cannot
help getting caught up in Escher's implied chain of levels, in which, for any one level,
there is always another level above it of greater "reality", and likewise, there is always a
level below, "more imaginary" than it is. This can be mind-boggling in itself. However,
what happens if the chain of levels is not linear, but forms a loop? What is real, then, and
what is fantasy? The genius of Escher was that he could not only concoct, but actually
portray, dozens of half-real, half-mythical worlds, worlds filled with Strange Loops,
which he seems to be inviting his viewers to enter.

                                           Gödel

In the examples we have seen of Strange Loops by Bach and Escher, there is a conflict
between the finite and the infinite, and hence a strong sense of paradox. Intuition senses
that there is something mathematical involved here. And indeed in our own century a
mathematical counterpart was discovered, with the most enormous repercussions. And,
just as the Bach and Escher loops appeal to very simple and ancient intuitions-a musical
scale, a staircase-so this discovery, by K. Gödel, of a Strange Loop in




Introduction: A Musico-Logical Offering                                                   23

                               FIGURE 9. Kurt Godel.




Introduction: A Musico-Logical Offering                24

mathematical systems has its origins in simple and ancient intuitions. In its absolutely
barest form, Godel's discovery involves the translation of an ancient paradox in
philosophy into mathematical terms. That paradox is the so-called Epimenides paradox,
or liar paradox. Epimenides was a Cretan who made one immortal statement: "All
Cretans are liars." A sharper version of the statement is simply "I am lying"; or, "This
statement is false". It is that last version which I will usually mean when I speak of the
Epimenides paradox. It is a statement which rudely violates the usually assumed
dichotomy of statements into true and false, because if you tentatively think it is true,
then it immediately backfires on you and makes you think it is false. But once you've
decided it is false, a similar backfiring returns you to the idea that it must be true. Try it!
   The Epimenides paradox is a one-step Strange Loop, like Escher's Print Gallery. But
how does it have to do with mathematics? That is what Godel discovered. His idea was to
use mathematical reasoning in exploring mathematical reasoning itself. This notion of
making mathematics "introspective" proved to be enormously powerful, and perhaps its
richest implication was the one Godel found: Godel's Incompleteness Theorem. What the
Theorem states and how it is proved are two different things. We shall discuss both in
quite some detail in this book. The Theorem can De likened to a pearl, and the method of
proof to an oyster. The pearl is prized for its luster and simplicity; the oyster is a complex
living beast whose innards give rise to this mysteriously simple gem.
   Godel's Theorem appears as Proposition VI in his 1931 paper "On Formally
Undecidable Propositions in Principia Mathematica and Related Systems I." It states:

    To every w-consistent recursive class K of formulae there correspond recursive
    class-signs r, such that neither v Gen r nor Neg (v Gen r) belongs to Fig (K) (where v
    is the free variable of r).

Actually, it was in German, and perhaps you feel that it might as well be in German
anyway. So here is a paraphrase in more normal English:

    All consistent axiomatic formulations of number theory
    include undecidable propositions.

This is the pearl.
   In this pearl it is hard to see a Strange Loop. That is because the Strange Loop is buried
in the oyster-the proof. The proof of Godel's Incompleteness Theorem hinges upon the
writing of a self-referential mathematical statement, in the same way as the Epimenides
paradox is a self-referential statement of language. But whereas it is very simple to talk
about language in language, it is not at all easy to see how a statement about numbers can
talk about itself. In fact, it took genius merely to connect the idea of self-referential
statements with number theory. Once Godel had the intuition that such a statement could
be created, he was over the major hurdle. The actual creation of the statement was the
working out of this one beautiful spark of intuition.




Introduction: A Musico-Logical Offering                                                     25

We shall examine the Godel construction quite carefully in Chapters to come, but so that
you are not left completely in the dark, I will sketch here, in a few strokes, the core of the
idea, hoping that what you see will trigger ideas in your mind. First of all, the difficulty
should be made absolutely clear. Mathematical statements-let us concentrate on number-
theoretical ones-are about properties of whole numbers. Whole numbers are not
statements, nor are their properties. A statement of number theory is not about a.
statement of number theory; it just is a statement of number theory. This is the problem;
but Godel realized that there was more here than meets the eye.
  Godel had the insight that a statement of number theory could be about a statement of
number theory (possibly even itself), if only numbers could somehow stand for
statements. The idea of a code, in other words, is at the heart of his construction. In the
Godel Code, usually called "Godel-numbering", numbers are made to stand for symbols
and sequences of symbols. That way, each statement of number theory, being a sequence
of specialized symbols, acquires a Godel number, something like a telephone number or a
license plate, by which it can be referred to. And this coding trick enables statements of
number theory to be understood on two different levels: as statements of number theory,
and also as statements about statements of number theory.
  Once Godel had invented this coding scheme, he had to work out in detail a way of
transporting the Epimenides paradox into a numbertheoretical formalism. His final
transplant of Epimenides did not say, "This statement of number theory is false", but
rather, "This statement of number theory does not have any proof". A great deal of
confusion can be caused by this, because people generally understand the notion of
"proof" rather vaguely. In fact, Godel's work was just part of a long attempt by
mathematicians to explicate for themselves what proofs are. The important thing to keep
in mind is that proofs are demonstrations within fixed systems of propositions. In the case
of Godel's work, the fixed system of numbertheoretical reasoning to which the word
"proof" refers is that of Principia Mathematica (P.M.), a giant opus by Bertrand Russell
and Alfred North Whitehead, published between 1910 and 1913. Therefore, the Godel
sentence G should more properly be written in English as:

      This statement of number theory does not have any proof in the system of Principia
      Mathematica.

Incidentally, this Godel sentence G is not Godel's Theorem-no more than the Epimenides
sentence is the observation that "The Epimenides sentence is a paradox." We can now
state what the effect of discovering G is. Whereas the Epimenides statement creates a
paradox since it is neither true nor false, the Godel sentence G is unprovable (inside
P.M.) but true. The grand conclusion\% That the system of Principia Mathematica is
"incomplete"-there are true statements of number theory which its methods of proof are
too weak to demonstrate.




Introduction: A Musico-Logical Offering                                                    26
